% Part: sets-functions-relations
% Chapter: sets
% Section: proofs-about-sets

\documentclass[../../../include/open-logic-section]{subfiles}

\begin{document}

\olfileid[fr]{sfr}{set}{prf}
\olsection{Démonstrations sur les ensembles}

\begin{editorial}
Pour un exposé général des techniques de démonstration,
voir le \olref[mth][prf][]{chap}.
On applique ici ces techniques à une loi d'absorption.
%
%
\end{editorial}

\begin{explain}
Les ensembles et les notations introduites jusqu'ici nous fournissent
des abréviations commodes pour décrire des ensembles et exprimer leurs
relations. Souvent, il faudra aussi démontrer des énoncés sur ces
relations. Si vous découvrez les démonstrations mathématiques,
commençons par un exemple simple. Nous allons montrer que, pour
tous ensembles $X$ et $Y$, on a toujours $X\cap(X\cup Y)=X$.
Comment démontrer une telle
égalité d'ensembles ? Rappelons que les ensembles sont déterminés
uniquement par leurs !!{element}s : deux ensembles sont égaux si et
seulement s'ils ont les mêmes !!{element}s. Nous devons donc établir
(a) que tout !!{element} de $X\cap(X\cup Y)$ appartient aussi à $X$,
et (b), réciproquement, que tout !!{element} de $X$ appartient aussi
à $X\cap(X\cup Y)$. Autrement dit, il faut montrer à la fois
(a) $X\cap(X\cup Y)\subseteq X$ et
(b) $X\subseteq X\cap(X\cup Y)$.

Pour démontrer un énoncé général comme « tout !!{element}~$z$ de
$X\cap(X\cup Y)$ appartient aussi à $X$ », on prend un~$z$
arbitraire, on suppose $z\in X\cap(X\cup Y)$, puis on déduit de cette
hypothèse que $z\in X$. Vous reconnaissez peut-être le schéma général
d'une preuve conditionnelle. Nous utiliserons aussi les définitions
des notions figurant dans l'hypothèse et dans la conclusion, ici
celles de~$\cap$ et de~$\cup$. Voici donc la première partie de la
démonstration :

\begin{quote}
(a) Nous voulons d'abord montrer que
$X\cap(X\cup Y)\subseteq X$, c'est-à-dire, par définition de
$\subseteq$, que $z\in X\cap(X\cup Y)$ entraîne $z\in X$ pour
tout~$z$. Supposons donc $z\in X\cap(X\cup Y)$. Un objet est
!!a{element} de l'intersection de deux ensembles si et seulement
s'il appartient à chacun d'eux. Nous pouvons donc conclure que
$z\in X$ et $z\in X\cup Y$. En particulier, $z\in X$, comme il
fallait le montrer.
\end{quote}

La première moitié de la démonstration est terminée. Remarquez que
la dernière étape utilise le fait que, si une conjonction
($z\in X$ et $z\in X\cup Y$) découle d'une hypothèse, chacun de ses
deux termes en découle également. Cette règle est l'« élimination de
la conjonction », ou $\land$Elim. Démontrons maintenant~(b) :

\begin{quote}
(b) Montrons que $X\subseteq X\cap(X\cup Y)$, c'est-à-dire, par
définition de $\subseteq$, que $z\in X$ entraîne
$z\in X\cap(X\cup Y)$ pour tout~$z$. Supposons $z\in X$. Pour
établir $z\in X\cap(X\cup Y)$, la définition de $\cap$ exige
(i) $z\in X$ et (ii) $z\in X\cup Y$. La condition~(i) est notre
hypothèse : il n'y a rien de plus à démontrer. Pour~(ii), rappelons
qu'un objet appartient à l'union de deux ensembles si et seulement
s'il appartient à au moins l'un des deux. Or $z\in X$ et
$X\cup Y$ est l'union de $X$ et de $Y$ ; ainsi, $z\in X\cup Y$.
Nous avons établi (i) $z\in X$ et (ii) $z\in X\cup Y$, d'où,
par définition de $\cap$, $z\in X\cap(X\cup Y)$.
\end{quote}

Cette démonstration est un peu longue, mais elle montre comment on
raisonne sur les ensembles et leurs relations. En général, on ne
détaille pas autant les choses et, notamment, on ne répète pas toutes
les définitions. Dans un texte plus avancé, la démonstration serait
beaucoup plus concise. Elle pourrait prendre la forme suivante.
\end{explain}

\begin{prop}[Absorption]
Pour tous ensembles $X$ et $Y$,
\[
X\cap(X\cup Y)=X.
\]
\end{prop}

\begin{proof}
(a) Supposons $z\in X\cap(X\cup Y)$. Alors $z\in X$, donc
$X\cap(X\cup Y)\subseteq X$.

(b) Supposons maintenant $z\in X$. Alors $z\in X\cup Y$ et, par
conséquent, $z\in X\cap(X\cup Y)$. Donc
$X\subseteq X\cap(X\cup Y)$.
\end{proof}

\begin{prob}
Démontrez en détail que $X\cup(X\cap Y)=X$, puis donnez-en une
démonstration abrégée. Indication : pour établir
$X\cup(X\cap Y)\subseteq X$, il faut raisonner par cas, c'est-à-dire
utiliser l'élimination de la disjonction ($\lor$Elim).
\end{prob}

\end{document}
